Sei $\ell$ die Länge und $m$ die Masse des Bretts, $M$ die Masse des Eimers und $x$ sein Abstand vom vorderen Ende. Die Gewichtskraft des Bretts greift in seiner Mitte an.

Nimm das vordere Ende als Drehachse: Dann fällt die Kraft $F_\mathrm{v}$ des vorderen Malers aus den Drehmomenten heraus. Die Kraft $F_\mathrm{h}$ des hinteren Malers dreht das Brett in die eine Richtung, die beiden Gewichtskräfte in die andere:
\begin{align}
 F_\mathrm{h}\cdot\ell &= m\,g\cdot\frac{\ell}{2} + M\,g\cdot(\ell-x)
\end{align}
Also
\begin{align}
 F_\mathrm{h} &= \FhF \\
   &= \frac{\mB\cdot\g\cdot\frac{\LB}{2}+\mE\cdot\g\cdot(\LB-\xE)}{\LB} \\
   &= \Fh \approx \result{\FhP{0}{3}}
\end{align}
Auch die Kräfte heben sich auf: Beide Maler zusammen tragen Brett und Eimer.
\begin{align}
 F_\mathrm{v} &= \FvF \\
   &= (\mB+\mE)\cdot\g-\Fh \\
   &= \Fv \approx \result{\FvP{0}{3}}
\end{align}
Der vordere Maler, näher beim Eimer, trägt mehr.
